\BookChapter{7}

\section[第7章\q 自由]{第7章\\ 自由}

{\kaishu

时间自由是财富能带给你的最大红利。

}

富有的最高级形式是，每天早上起床后你都可以说：“今天我能做我想做的任何事。”

人们都想变得更富有，希望财富能带来更多幸福。幸福是一个复杂的话题，因为每个人的幸福观都不同，但如果幸福的分数有一个公分母——一种普遍的快乐源泉——那就是对生活的全面掌控。

在自己喜欢的任何时候和自己喜欢的对象做想做的事，而且想做多久就做多久，这样的自由是极其珍贵的，而这就是金钱能带给我们的最大红利。

安格斯·坎贝尔（Angus Campbell）生于1910年，是密歇根大学的心理学家。在他进行研究的那个年代，心理学的绝大多数关注点落在那些影响人们行使正常功能的心理障碍上，比如抑郁、焦虑和精神分裂症。

坎贝尔想知道的却是什么会带给人幸福。他在1981年出版的图书《美国人的幸福感》（\textit{The Sense of Well-being in America}）中开篇便指出，人们普遍比很多心理学家认为的要快乐得多，但一些人的幸福感显然高于其他人。我们不能按收入、地区或受教育程度简单地判断谁比谁更快乐，因为每个类别中都有很多人长期感到不快乐。

幸福最强大的公分母很简单。坎贝尔总结道：

{\kaishu

与我们考虑过的任何客观生活条件相比，对自己的生活拥有强烈的掌控感是预测幸福这种积极感受的一项更可靠的指标。

}

不是工资多少，不是房子大小，也不是工作好坏，而是对自己想做什么、什么时候做、和谁一起做拥有掌控能力。这是生活中决定幸福感的通用变量。

金钱最大的内在价值是它能赋予你掌控自己时间的能力——这句话没有任何夸张的成分。你之所以能从财富积累过程中逐步获得一定程度的独立和自主，是因为你积累的财富能让你对可以做什么以及什么时候去做有更强的控制力。

拥有一小笔财富意味着生病时就算请几天假也不至于入不敷出。对没有这种能力的人来说，拥有这种能力会给生活带来巨大的进步。

拥有更多财富则意味着在失业后可以从容地等待更好的职业机会，而不必急于抓住遇到的第一根救命稻草。这种能力可以改变一个人的生活。

拥有6个月的应急存款意味着在面对老板时不至于如履薄冰，因为你知道即使需要时间另谋出路，你的生活也不会陷入窘境。

拥有更多财富则意味着可以选择一份待遇不高但时间灵活的，或是通勤时间比较短的工作。这也意味着在遇到突发疾病时不会出现对付不起高额医疗账单的多余担心。

除此之外，你还可以自由选择退休的时间，而不必一把年纪了还在为生计奔波。

财富能提供你更多的时间和选项。它赋予你生活的好处，没有什么奢侈品比得上。

我在整个大学期间心心念念地想去投资银行工作。原因只有一个：这份工作能赚很多钱。这是我唯一的动力，我也100\%相信只要我做了这份工作，我一定会更快乐。大三时，我在洛杉矶的一家投资银行获得了一个暑期实习的机会。我当时觉得自己拿到了职业生涯的入场券，而这就是我梦寐以求的全部了。

在上班的第一天，我就明白了为什么投资银行家们会赚到那么多钱：他们工作的时长和辛苦程度超过了人类能承受的限度。事实上，大多数人都受不了这样的工作。在半夜前回家都是一种奢侈的事。这一行流传着这样一种说法：“如果你周六不来上班，周日就不用来了。”这份工作考验智慧，报酬丰厚，能让你体现自己的价值，但你在醒着的每一秒都要为老板提出的各种要求疲于奔命。于是，这份工作成了我人生中最悲惨的经历之一。实习期有4个月，而我只坚持了1个月。

最令我介怀的一点是，我喜欢这份工作，也想努力做好它，但是做一份自己喜欢却无法掌控时间的事和做自己讨厌的事没什么区别。

这种感觉在心理学中有一个专门的名称——抗拒（reactance）。宾夕法尼亚大学的营销学教授乔纳·伯杰（Jonah Berger）总结道：

{\kaishu

人们喜欢一切尽在自己掌握的感觉，就好像坐在生活的驾驶席上。当他们被要求做一些事时，他们会觉得自己没有话语权，没有选择权，选择是别人替他们做出的。所以，他们会拒绝做这件事或故意去做其他事，即使他们本来可能很乐意做这件事。\footnote{J. Berger,“How to Change Anyone’s Mind,” The Wall Street Journal (February 21, 2020).}

}

如果你认可这句话，你就会意识到，能让你随心所欲地在任何时间、任何地点，与任何人一起做任何你想做的事，而且想做多久做多久的理财方式，才是财富最令人惊叹的回报。

德雷克·西弗斯（Derek Sivers）是一名成功的企业家。他曾写到一位朋友让他讲讲他是如何致富的：

{\kaishu

我在曼哈顿市中心有一份工作，年收入是2万美元——这大概是最低工资水平了……我从不在饭馆吃饭，也从不打车。我每月的生活开支大约是1,000美元，而我的月薪是1,800美元。我工作了两年，存了1.2万美元。当时我22岁。

有了这1.2万，我就可以辞掉工作，做一名全职音乐家了。我知道每个月我都可以接到一些演出邀约来支付生活费。这样一来，我就自由了。一个月后我辞掉了工作，之后再也没有给别人打过工。

在我给朋友讲完故事后，他让我继续讲下去。我说没了，这就是全部故事了。他说：“那你卖掉公司的事呢？”

我说，不，那件事没对我的生活造成太大影响，只意味着我的银行存款增加了而已。我生活中真正的影响是在22岁时发生的。\footnote{D. Sivers,“How I got rich on the other hand,” sivers.org (October 30, 2019).}

}

美国是世界历史上最富有的国家，但是几乎没什么证据可以证明美国人现在比20世纪50年代时更幸福，而那时美国人的财富和收入比现在要低得多——即使是以去掉通货膨胀因素后的中位数水平来考量。根据盖洛普（Gallup）民意测验中心在2019年对140个国家的15万人进行的调查，大约有45\%的美国人表示他们前一天有过“非常忧虑”的情绪，\footnote{N. Chokshi,“Americans Are Among the Most Stressed People in the World, Poll Finds,” The New York Times (April 25, 2019).}而全球平均水平是39\%。55\%的美国人表示他们前一天有过“压力巨大”的感受，而在世界其他地区，这样的回答只占全体的35\%。

这种情况的部分原因是，美国人把很多钱花在买更大、更好的商品上，与此同时也就放弃了对自己时间的更多控制。就算在理想情况下，前者的好处也会被后者的害处抵消。

1955年，美国家庭收入的中位数是2.9万美元（经通货膨胀调整后），\footnote{Russell Sage Foundation—Chartbook of Social Inequality.}到2019年，这个数值变成了6.2万美元。拥有这些财富的美国人所过的生活是20世纪50年代的美国人无法想象的，即使是对当时中等收入的家庭来说。美国中等收入家庭的房屋面积中位数从1950年的91平方米涨到了2018年的226平方米。美国如今的新建住房中卫生间的数量都比居住人数多。美国人的汽车更快、更节能了，电视也更便宜、画质更好了。

但从另一个方面看，在这个时代发生的事看上去并不像进步。这和现在许多美国人从事的工作有很大关系。

约翰·D.洛克菲勒是有史以来最成功的商人之一。他还是个隐士，大部分时间都在独处。他沉默寡言，尽量减少交际，就算你引起了他的注意，他也只会保持沉默。

成为洛克菲勒心腹的一名炼油工人曾说：“他让每个人都发言，而他自己却靠在椅背上，不说话。”

当有人问起他为什么在会议中保持沉默时，洛克菲勒经常会引用一首诗来回答：

{\kaishu

一只智慧的老猫头鹰住在橡树上，

它看得越多，说得就越少，

它说得越少，听得就越多，

为什么我们不能像这只智慧的老猫头鹰一样呢？

}

洛克菲勒是个脾气古怪的家伙，但他的回答却适用于今天数以千万计的上班族。

洛克菲勒的工作内容不是钻井、装载火车或者搬运石油，而是在思考后做出正确的决定。洛克菲勒的产品——他需要交付的东西——并不是他用双手做的，甚至也不是他用嘴说的，而是他在脑中思考出来的，而那才是他花费大部分时间和精力的地方。尽管在很多人看来，他每天大多时候只是静静地坐着，看起来无所事事，但实际上他一直在脑中工作，在思考问题。

这种情况在他那个时代是很少见的。在洛克菲勒生活的美国，几乎所有工作都需要用双手去做。据经济学家罗伯特·戈登（Robert Gordon）统计，1870年，美国46\%的工作是农业相关的，35\%是手工业或制造业相关的。很少有哪个职业需要依靠工作者的大脑。你不用思考，只需要不断地用手干活。你的工作是看得见、摸得着的。

今天，一切大不一样了。

如今，有38\%的工作由“经理”“高级职员”和“专业人士”来完成，都是决策类工作；另有41\%是服务类工作，既需要动手，也常常需要动脑筋。

我们当中越来越多人的工作类型趋近洛克菲勒当初的工作类型，而不像20世纪50年代制造业工人们的工作类型。这意味着我们每天离开工位、打卡下班后，工作却并未结束。我们还在脑海中持续地劳动，就好像工作永远不会结束一样。

如果你的工作是制造汽车，在离开装配线后，几乎没有什么是你可以做的了。一旦你放下工作，把工具留在工厂里，你的工作就完全停止了。但如果你的工作是营销策划——一个基于脑力劳动和决策的工作——你的工具就是你的大脑，而这个工具永远不会离开你。这样一来，在通勤的路上，在做晚餐时，在哄孩子睡觉时，在凌晨3点因为压力醒来时，你或许都在想着你的项目。你虽然正常上下班，但感觉好像每时每刻都在工作。

《大西洋月刊》（\textit{The Atlantic}）的撰稿人德雷克·汤普森（Derek Thompson）曾这样描述：

{\kaishu

如果21世纪的操作设备是一个便携式装置，那么这就意味着现代工厂压根不是一个作为实体的场所，而是一天的时间本身。计算机时代把办公工具从办公室中解放了出来。大多数脑力工作者的笔记本电脑和智能手机是便携式的全能媒体制造机器。从理论上讲，无论是下午2点在办公室，凌晨2点在东京的Wework\footnote{总部位于美国纽约的共享办公企业，提供共享办公空间。}写字楼，还是午夜坐在沙发上，你都能以同样的效率工作。\footnote{D. Thompson,“Why White-Collar Workers Spend All Day at the Office,” The Atlantic (December 4, 2019).}

}

这就意味着，与前几代人相比，我们对时间的控制力降低了。正因为控制时间是影响幸福感的关键因素，所以我们无须对尽管现在的我们更富有了，但我们没有感到更快乐这一事实感到惊讶。

那么我们该怎么做呢？

这个问题并不好解决，因为每个人的情况不同。解决的第一步是，我们需要明确什么会让每个人都快乐，什么不会。

在《有关生活的30条经验》（\textit{30 Lessons for Living}）中，老年学家卡尔·皮尔默（Karl Pillemer）采访了1,000名美国老年人，以总结他们从几十年的生活中得到的最重要的教训。他写道：

{\kaishu

没有人——1,000个人中没有1个——表示如果你想变得快乐，你就该拼命地工作赚钱，然后买自己想要的东西。

没有人——没有1个——表示你至少要和周围的人一样富有，或者要比周围的人更富有才算真正成功。

没有人——没有1个——表示你应该根据你期望获得的收入水平来选择你的工作。

}

这些受采访的老人看重的是获得真诚的友谊，参与比个体的存在更伟大的事业，把大量闲暇好时光花在跟子女相处上。“比起你的钱（或者你的钱能买到的东西），你的孩子们更想要的是和你在一起的体验。”皮尔默写道。

那些历经世事沧桑的人能给你的忠告是：时间自由是财富能带给你的最大红利。

现在，让我们用短短的一章来介绍财富能带给你的一种最微不足道的红利。
